\documentclass[10pt,a4paper]{amsart}
\usepackage{amsmath,amssymb,amsthm,graphicx}
% Retained minimal preamble for the Stacks Project chapter more-algebra.tex (More on Algebra).
% The chapter file pinned-original-source.tex carries only \input{preamble} on its first line, so no
% preamble text lives inside the chapter itself.  The Stacks Project uses one single shared preamble for
% every chapter; the snapshot stages it identically as sources/stacks-chow/preamble.tex and
% sources/stacks-etale-cohomology/preamble.tex.  The block below is copied mechanically, line for line,
% from sources/stacks-chow/preamble.tex lines 14-18 (the Xy-pic setup) and 154-217 (the native theorem
% declarations together with the mathematical macros).  Nothing was added, removed or redefined; the
% wrapper class and the amsmath/amssymb/amsthm/graphicx base packages are supplied by the assembler.
% Chapter: more-algebra.tex -- More on Algebra.
% For commutative diagrams we use Xy-pic
\usepackage[all]{xy}

% We use 2cell for 2-commutative diagrams.
\xyoption{2cell}

% Theorem environments.
%
\theoremstyle{plain}
\newtheorem{theorem}[subsection]{Theorem}
\newtheorem{proposition}[subsection]{Proposition}
\newtheorem{lemma}[subsection]{Lemma}

\theoremstyle{definition}
\newtheorem{definition}[subsection]{Definition}
\newtheorem{example}[subsection]{Example}
\newtheorem{exercise}[subsection]{Exercise}
\newtheorem{situation}[subsection]{Situation}

\theoremstyle{remark}
\newtheorem{remark}[subsection]{Remark}
\newtheorem{remarks}[subsection]{Remarks}

\numberwithin{equation}{subsection}

% Macros
%
\def\lim{\mathop{\mathrm{lim}}\nolimits}
\def\colim{\mathop{\mathrm{colim}}\nolimits}
\def\Spec{\mathop{\mathrm{Spec}}}
\def\Hom{\mathop{\mathrm{Hom}}\nolimits}
\def\Ext{\mathop{\mathrm{Ext}}\nolimits}
\def\SheafHom{\mathop{\mathcal{H}\!\mathit{om}}\nolimits}
\def\SheafExt{\mathop{\mathcal{E}\!\mathit{xt}}\nolimits}
\def\Sch{\mathit{Sch}}
\def\Mor{\mathop{\mathrm{Mor}}\nolimits}
\def\Ob{\mathop{\mathrm{Ob}}\nolimits}
\def\Sh{\mathop{\mathit{Sh}}\nolimits}
\def\NL{\mathop{N\!L}\nolimits}
\def\CH{\mathop{\mathrm{CH}}\nolimits}
\def\proetale{{pro\text{-}\acute{e}tale}}
\def\etale{{\acute{e}tale}}
\def\QCoh{\mathit{QCoh}}
\def\Ker{\mathop{\mathrm{Ker}}}
\def\Im{\mathop{\mathrm{Im}}}
\def\Coker{\mathop{\mathrm{Coker}}}
\def\Coim{\mathop{\mathrm{Coim}}}

% Boxtimes
%
\DeclareMathSymbol{\boxtimes}{\mathbin}{AMSa}{"02}

%
% Macros for moduli stacks/spaces
%
\def\QCohstack{\mathcal{QC}\!\mathit{oh}}
\def\Cohstack{\mathcal{C}\!\mathit{oh}}
\def\Spacesstack{\mathcal{S}\!\mathit{paces}}
\def\Quotfunctor{\mathrm{Quot}}
\def\Hilbfunctor{\mathrm{Hilb}}
\def\Curvesstack{\mathcal{C}\!\mathit{urves}}
\def\Polarizedstack{\mathcal{P}\!\mathit{olarized}}
\def\Complexesstack{\mathcal{C}\!\mathit{omplexes}}
% \Pic is the operator that assigns to X its picard group, usage \Pic(X)
% \Picardstack_{X/B} denotes the Picard stack of X over B
% \Picardfunctor_{X/B} denotes the Picard functor of X over B
\def\Pic{\mathop{\mathrm{Pic}}\nolimits}
\def\Picardstack{\mathcal{P}\!\mathit{ic}}
\def\Picardfunctor{\mathrm{Pic}}
\def\Deformationcategory{\mathcal{D}\!\mathit{ef}}

% --- appended 2026-10-12 (Chief): the project own definitions of the reference/slogan/history
% display environments, copied mechanically from the shared Stacks preamble lines 8-12
% (\usepackage{verbatim} plus the three \newenvironment lines, which the chapter preamble keeps
% verbatim from the project preamble). Nothing added, removed or redefined: these environments
% carry no mathematical content and the project itself typesets them as comments. Without these
% three lines a counted statement whose own text uses them fails with Environment slogan undefined.
\usepackage{verbatim}
\newenvironment{reference}{\comment}{\endcomment}
\newenvironment{slogan}{\comment}{\endcomment}
\newenvironment{history}{\comment}{\endcomment}
\usepackage{cleveref}
\crefname{originallemma}{original lemma}{original lemmas}
\Crefname{originallemma}{Original lemma}{Original lemmas}
\makeatletter
\newlabel{lemma-ses}{{36678}{source}{}{sourcepointer-1}{}}
\newlabel{lemma-ses@cref}{{[originallemma][0][]36678}{[1][0][]source}{}{}{}}
\makeatother
\title{For a flat local homomorphism the determinant of a base changed endomorphism is the m-th power of the original determinant}
\author{The Stacks Project Authors}
\date{stacks-project:more-algebra.tex:local-snapshot-20261012; GNU FDL 1.2 or later}
\begin{document}
\maketitle
\setcounter{subsection}{1}
\noindent\textit{Editorial note.} All delivered mathematical text is the author's own, copied from the pinned Stacks Project chapter file more-algebra.tex; no statement or proof step was rewritten, reordered or removed.\par\medskip
\noindent\textit{Editorial note.} The chapter file pinned-original-source.tex carries only \textbackslash{}input\{preamble\} on its first line, so no preamble text lives inside the chapter. The retained file sources/stacks-more-algebra/retained-preamble.tex is a mechanical, line-for-line copy of the project's single shared preamble as staged in the snapshot; nothing was added, removed or redefined and no package was added.\par\medskip
\noindent\textit{Editorial note.} The source URL is the source-level Stacks Project URL, as required when more than one card is produced from one source; the per-claim tag is carried by claim\_key.\par\medskip
\noindent\textit{Editorial note.} The statement refers to the exponent m defined in its own first display and to the object (M, varphi) 'as above', both delivered verbatim. The three references inside the counted proof span are recorded as non-load-bearing navigation pointers.\par\medskip
\noindent\textit{Allowed standard prerequisites (editorial).}\begin{itemize}
\item The determinant det\_kappa(varphi) of an endomorphism of a finite length module over the residue field, and its multiplicativity along short exact sequences. invoked through the earlier lemma at proof lines 36735 and 36755. Is the invariant being base changed and the tool that reduces to simple objects.
\item Flatness of the local homomorphism, giving base change of short exact sequences. used at proof lines 36733-36734. Is what makes the reduction to simple objects legitimate.
\item Compatibility of determinants of linear maps with field extensions. used at proof lines 36763-36764. Identifies the determinant over kappa' of the base changed map on the simple quotient with the m-th power of the determinant over kappa.
\end{itemize}\medskip
\ifdefined\hypertarget\hypertarget{sourcepointer-1}{}\fi
\noindent\textit{Source reference.} Original lemma 36678: sources/stacks-more-algebra/pinned-original-source.tex, source lines 36678--36678. Same-chapter navigation pointer; the statement it names is standard for the section and the argument that consumes it stays inside the counted proof span.\par\medskip
\begin{lemma}
\label{lemma-flat-base-change-det}
Let $(R, \mathfrak m, \kappa) \to (R', \mathfrak m', \kappa')$
be a flat local homomorphism of local rings such that
$m = \text{length}_{R'}(R'/\mathfrak mR') < \infty$.
For any $(M, \varphi)$ as above, the element
$\det_\kappa(\varphi)^m$ maps to
$\det_{\kappa'}(\varphi \otimes 1 : M \otimes_R R' \to M \otimes_R R')$
in $\kappa'$.
\end{lemma}

\begin{proof}
The flatness of $R \to R'$ assures us that short exact sequences
as in Lemma \ref{lemma-ses} base change to short exact sequences
over $R'$. Hence by the multiplicativity of Lemma \ref{lemma-ses}
we may assume that $(M, \varphi)$ is a simple object of our category
(see introduction to this section). In the simple case $M$ is annihilated
by $\mathfrak m$. Choose a filtration
$$
0 \subset I_1 \subset I_2 \subset \ldots \subset I_{m - 1} \subset
R'/\mathfrak mR'
$$
whose successive quotients are isomorphic to $\kappa'$ as $R'$-modules.
Then we obtain the filtration
$$
0 \subset
M \otimes_\kappa I_1 \subset
M \otimes_\kappa I_2 \subset
\ldots \subset
M \otimes_\kappa I_{m - 1} \subset
M \otimes_\kappa R'/\mathfrak mR' = M \otimes_R R'
$$
whose successive quotients are isomorphic to $M \otimes_\kappa \kappa'$.
Also, these submodules are invariant under $\varphi \otimes 1$.
By Lemma \ref{lemma-ses} we find
$$
\det\nolimits_{\kappa'}(\varphi \otimes 1 : M \otimes_R R' \to M \otimes_R R')
=
\det\nolimits_{\kappa'}(\varphi \otimes 1 :
M \otimes_\kappa \kappa' \to M \otimes_\kappa \kappa')^m =
\det\nolimits_\kappa(\varphi)^m
$$
The last equality holds by the compatibility of determinants
of linear maps with field extensions. This proves the lemma.
\end{proof}

\end{document}
